% Part: first-order-logic
% Chapter: syntax-and-semantics
% Section: first-order-languages

\documentclass[../../../include/open-logic-section]{subfiles}

\begin{document}

\olfileid{fol}{syn}{fol}

\olsection{প্রথম-ক্রমের ভাষা}


প্রথম-ক্রম যুক্তিবিদ্যার রাশিগুলি একটি মৌলিক শব্দভাণ্ডার থেকে গড়ে ওঠে;
এতে থাকে \emph{!!{variable}s}, \emph{!!{constant}s},
\emph{!!{predicate}s} এবং কখনও কখনও \emph{!!{function}s}। এগুলির
সঙ্গে যৌক্তিক সংযোজক, পরিমাণসূচক এবং বন্ধনী ও কমার মতো বিরামচিহ্ন
ব্যবহার করে \emph{পদ} ও \emph{!!{formula}s} গঠন করা হয়।

\begin{explain}
অনানুষ্ঠানিকভাবে, !!{predicate}s হলো ধর্ম ও সম্বন্ধের নাম,
!!{constant}s হলো একক বস্তুর নাম এবং !!{function}s হলো প্রতিচিত্রের
নাম। !!{identity}~$\eq$ ছাড়া এগুলিই \emph{অ-যৌক্তিক সংকেত}, আর
সমষ্টিগতভাবে এগুলিই একটি ভাষা গঠন করে। যেকোনো প্রথম-ক্রমের ভাষা~$\Lang
L$ তার অ-যৌক্তিক সংকেতগুলি দ্বারা নির্ধারিত। সবচেয়ে সাধারণ ক্ষেত্রে
$\Lang L$-এ প্রত্যেক প্রকারের অসীম সংখ্যক সংকেত থাকে।
\end{explain}

সাধারণ ক্ষেত্রে প্রথম-ক্রম যুক্তিবিদ্যায় আমরা নিচের সংকেতগুলি ব্যবহার
করি:

\begin{enumerate}
\item যৌক্তিক সংকেত
\begin{enumerate}
\item যৌক্তিক সংযোজক:
  \startycommalist
  \iftag{prvNot}{\ycomma $\lnot$ (নঞর্থকরণ)}{}%
  \iftag{prvAnd}{\ycomma $\land$ (সংযোজন)}{}%
  \iftag{prvOr}{\ycomma $\lor$ (বিয়োজন)}{}%
  \iftag{prvIf}{\ycomma $\lif$ (!!{conditional})}{}%
  \iftag{prvIff}{\ycomma $\liff$ (!!{biconditional})}{}%
  \iftag{prvAll}{\ycomma $\lforall$ (সার্বিক পরিমাণসূচক)}{}%
  \iftag{prvEx}{\ycomma $\lexists$ (অস্তিত্বসূচক পরিমাণসূচক)}{}।
\tagitem{prvFalse}{!!{falsity}-এর বচনমূলক ধ্রুবক~$\lfalse$।}{}
\tagitem{prvTrue}{!!{truth}-এর বচনমূলক ধ্রুবক~$\ltrue$।}{}
\item দ্বি-স্থানীয় !!{identity}~$\eq$।
\item !!{variable}s-এর একটি !!{denumerable}s সেট: $\Obj v_0$, $\Obj v_1$, $\Obj
  v_2$, \dots
\end{enumerate}
\item প্রথম-ক্রম যুক্তিবিদ্যার \emph{মানক ভাষা} গঠনকারী অ-যৌক্তিক
  সংকেত
\begin{enumerate}
\item প্রতি $n>0$-এর জন্য $n$-স্থানীয় !!{predicate}s-এর একটি
  !!{denumerable}s সেট: $\Obj A^n_0$, $\Obj A^n_1$, $\Obj A^n_2$, \dots
\item !!{constant}s-এর একটি !!{denumerable}s সেট: $\Obj c_0$, $\Obj c_1$, $\Obj
  c_2$, \dots।
\item প্রতি $n>0$-এর জন্য $n$-স্থানীয় !!{function}s-এর একটি
  !!{denumerable}s সেট:
  $\Obj f^n_0$, $\Obj f^n_1$, $\Obj f^n_2$, \dots
\end{enumerate}
\item বিরামচিহ্ন: (, ) এবং কমা।
\end{enumerate}

আমাদের অধিকাংশ সংজ্ঞা ও ফলাফল প্রথম-ক্রম যুক্তিবিদ্যার পূর্ণ মানক ভাষার
জন্য বিবৃত হবে। তবে প্রয়োগের ধরন অনুযায়ী ভাষাটিকে অল্প কয়েকটি
!!{predicate}s, !!{constant}s ও !!{function}s-এও সীমিত করতে পারি।

\begin{ex}
পাটিগণিতের ভাষা~$\Lang L_A$-তে আছে একটিমাত্র দ্বি-স্থানীয়
!!{predicate}~$<$, একটিমাত্র !!{constant}~$\Obj 0$, একটি এক-স্থানীয়
!!{function}~$\prime$ এবং দুটি দ্বি-স্থানীয় !!{function}s~$+$ ও~$\times$।
\end{ex}

\begin{ex}
সেটতত্ত্বের ভাষা~$\Lang L_Z$-তে কেবল একটিমাত্র দ্বি-স্থানীয়
!!{predicate}~$\in$ আছে।
\end{ex}

\begin{ex}
ক্রমের ভাষা~$\Lang L_\le$-তে কেবল দ্বি-স্থানীয়
!!{predicate}~$\le$ আছে।
\end{ex}

আবারও বলি, এগুলি প্রথামাত্র: বিধিবদ্ধভাবে এগুলি কেবল বিকল্প নাম। যেমন,
$<$, $\in$ ও $\le$ হলো $\Obj A^2_0$-এর বিকল্প নাম, $\Obj 0$ হলো
$\Obj c_0$-এর, $\prime$ হলো $\Obj f^1_0$-এর, $+$ হলো $\Obj f^2_0$-এর
এবং $\times$ হলো $\Obj f^2_1$-এর বিকল্প নাম।

\iftag{defNot,defOr,defAnd,defIf,defIff,defTrue,defFalse,defEx,defAll}{%
উপরে প্রবর্তিত মৌলিক সংযোজক ও
\iftag{notprvEx,notprvAll}{পরিমাণসূচকটির}{পরিমাণসূচকগুলির} পাশাপাশি
আমরা নিচের \emph{সংজ্ঞায়িত} সংকেতগুলিও ব্যবহার করি:
\startycommalist
  \iftag{defNot}{\ycomma $\lnot$ (নঞর্থকরণ)}{}%
  \iftag{defAnd}{\ycomma $\land$ (সংযোজন)}{}%
  \iftag{defOr}{\ycomma $\lor$ (বিয়োজন)}{}%
  \iftag{defIf}{\ycomma $\lif$ (!!{conditional})}{}%
  \iftag{defIff}{\ycomma $\liff$ (!!{biconditional})}{}%
  \iftag{defAll}{\ycomma $\lforall$ (সার্বিক পরিমাণসূচক)}{}%
  \iftag{defEx}{\ycomma $\lexists$ (অস্তিত্বসূচক পরিমাণসূচক)}{}%
  \iftag{defFalse}{\ycomma !!{falsity}~$\lfalse$}{}%
  \iftag{defTrue}{\ycomma !!{truth}~$\ltrue$}}{}।

\begin{tagblock}{defNot,defOr,defAnd,defIf,defIff,defTrue,defFalse,defEx,defAll}
\begin{explain}
সংজ্ঞায়িত সংকেত বিধিবদ্ধভাবে ভাষার অংশ নয়; অনানুষ্ঠানিক সংক্ষিপ্ত রূপ
হিসেবে তাকে প্রবর্তন করা হয়। কেবল মৌলিক সংকেত ব্যবহার করলে যে সূত্রগুলি
বেশ দীর্ঘ হয়ে যেত, এর সাহায্যে আমরা সেগুলিকে সংক্ষেপে লিখতে পারি। এটি
অবশ্যই একটি সুবিধা। তবে আরও বড় সুবিধা হলো, প্রমাণগুলি ছোট হয়। কোনো
সংকেত মৌলিক হলে প্রমাণে তাকে আলাদাভাবে বিবেচনা করতে হয়। তাই মৌলিক
সংকেত যত বেশি, প্রমাণও তত দীর্ঘ।
\end{explain}
\end{tagblock}

% Alternate symbols

\begin{intro}
উপরে ব্যবহৃত পরিভাষা ও সংকেতের বদলে অন্য পরিভাষা ও সংকেতের সঙ্গে আপনার
পরিচয় থাকতে পারে। যুক্তিবিদ্যার বইগুলি (এবং শিক্ষকেরা) “নঞর্থকরণ”-এর
জন্য সাধারণত $\sim$, $\neg$ বা~!\ এবং “সংযোজন”-এর জন্য $\wedge$,
$\cdot$ বা $\&$ ব্যবহার করেন। “শর্তবচন” বা “নিহিতকরণ”-এর প্রচলিত
সংকেত হলো $\rightarrow$, $\Rightarrow$ ও $\supset$।
\iftag{prvIff,defIff}{“দ্বিশর্তবচন”, “দ্বি-নিহিতকরণ” বা “(বস্তুগত)
  সমতুল্যতা”-র সংকেত হলো $\leftrightarrow$, $\Leftrightarrow$ ও
  $\equiv$।}{}
\iftag{prvFalse,defFalse}{$\lfalse$ সংকেতটিকে “মিথ্যাত্ব”, “ফলসাম”,
  “অসংগতি” বা “বটম”—নানা নামে ডাকা হয়।}{}
\iftag{prvTrue,defTrue}{$\ltrue$ সংকেতটিকে “সত্যতা”, “ভেরাম” বা
  “টপ”—নানা নামে ডাকা হয়।}{}

!!{constant}s-এর (কখনও যাদের নাম বলা হয়) জন্য লাতিন বর্ণমালার গোড়ার
ছোট হাতের অক্ষর (যেমন $a$, $b$, $c$) এবং !!{variable}s-এর জন্য শেষের
ছোট হাতের অক্ষর (যেমন $x$, $y$, $z$) ব্যবহার করাই প্রথা। পরিমাণসূচক
!!{variable}s-এর সঙ্গে যুক্ত হয়, যেমন $x$; সার্বিক পরিমাণসূচকের
সংকেতরূপের মধ্যে আছে $\forall x$, $(\forall x)$, $(x)$, $\Pi x$,
$\bigwedge_x$, আর অস্তিত্বসূচক পরিমাণসূচকের মধ্যে আছে $\exists x$,
$(\exists x)$, $(Ex)$, $\Sigma x$, $\bigvee_x$।
\end{intro}

\begin{explain}
আমরা সব বচনমূলক অপারেটর এবং উভয় পরিমাণসূচককে ভাষার মৌলিক সংকেত
হিসেবে ধরতে পারি। আবার মৌলিক সংকেতের একটি ছোটতর ভাণ্ডার বেছে নিয়ে
অন্য !!{operator}s-কে সংজ্ঞায়িত সংকেতও ধরতে পারি। বুলীয় অপারেটরের
“সত্য-ফলনগতভাবে পূর্ণ” সেটের মধ্যে আছে $\{ \lnot, \lor \}$,
$\{ \lnot, \land \}$ এবং $\{ \lnot, \lif\}$—এগুলির যেকোনোটিকে
যেকোনো একটি পরিমাণসূচকের সঙ্গে মিলিয়ে অভিব্যক্তিগতভাবে পূর্ণ
প্রথম-ক্রমের ভাষা পাওয়া যায়।

আরও দুটি !!{operator}s-এর সঙ্গে আপনার পরিচয় থাকতে পারে: হেনরি শেফারের
নামানুসারে শেফার স্ট্রোক~$|$ এবং পিয়ার্সের তির~$\downarrow$, যা
কোয়াইনের ড্যাগার নামেও পরিচিত। এগুলিকে যথাক্রমে “ন্যান্ড” ও “নর”-এর
প্রচলিত অর্থ দিলে প্রত্যেকটি একাই সত্য-ফলনগতভাবে পূর্ণ।
\end{explain}

\end{document}
